\documentclass[11pt,a4paper]{article}
\usepackage[utf8]{inputenc}
\usepackage[T1]{fontenc}
\usepackage[margin=2.2cm]{geometry}
\usepackage{amsmath,amssymb}
\usepackage{enumitem}
\usepackage{hyperref}

\setlength{\parindent}{0pt}
\setlength{\parskip}{0.75em}

\title{MA087G Mathematical Statistics\\Beginner-Friendly Full Solution Guide\\2025-01-10}
\author{}
\date{}

\begin{document}
\maketitle

\section*{Source notes before solving}
The official exam PDF was image-only. Normal text extraction found no text, but the rendered page images were readable.

Per request, this folder does not include an extraction log or mistake bank.

\section*{Uppgift 1}

\subsection*{What this question is really asking}
This is a three-set counting problem. There are 30000 students, and students can be active in areas \(A\), \(B\), and \(C\).

We must find the probability that a randomly chosen student is active in:
\begin{itemize}[leftmargin=*]
    \item at least one area,
    \item no area,
    \item exactly one area.
\end{itemize}

\subsection*{How to recognize this type of question}
The clues are:
\begin{itemize}[leftmargin=*]
    \item The question gives counts for \(A\), \(B\), and \(C\).
    \item It gives overlaps like both \(A\) and \(B\).
    \item It gives the number active in all three.
\end{itemize}
This is a Venn diagram/inclusion-exclusion problem.

\subsection*{Tiny theory needed}
For three events:
\[
|A\cup B\cup C|
=|A|+|B|+|C|-|A\cap B|-|A\cap C|-|B\cap C|+|A\cap B\cap C|.
\]
The vertical bars mean ``number of people in''.

We subtract pair-overlaps because they were counted twice. We add the triple-overlap back because it was subtracted too many times.

\subsection*{Formula and what everything means}
\[
P(\text{event})=\frac{\text{number of favorable students}}{\text{total number of students}}.
\]
Here the total is \(30000\).

For exactly one area:
\[
\text{only }A=|A|-|A\cap B|-|A\cap C|+|A\cap B\cap C|.
\]
We add the triple back because it was removed twice.

\subsection*{Step-by-step solution}
Given:
\[
|A|=12000,\quad |B|=8000,\quad |C|=6000.
\]
\[
|A\cap B|=7000,\quad |A\cap C|=4500,\quad |B\cap C|=1000.
\]
\[
|A\cap B\cap C|=500.
\]

\textbf{Part a: At least one area.}

At least one means \(A\cup B\cup C\). Use inclusion-exclusion:
\begin{align*}
|A\cup B\cup C|
&=12000+8000+6000\\
&\quad -7000-4500-1000+500\\
&=26000-12500+500\\
&=14000.
\end{align*}
So:
\begin{align*}
P(\text{at least one})
&=\frac{14000}{30000}\\
&=\frac{7}{15}\\
&\approx 0.467.
\end{align*}

\textbf{Part b: No area.}

No area is the complement of at least one area:
\begin{align*}
P(\text{no area})
&=1-P(\text{at least one})\\
&=1-\frac{7}{15}\\
&=\frac{8}{15}\\
&\approx 0.533.
\end{align*}

\textbf{Part c: Exactly one area.}

Exactly one means only \(A\), only \(B\), or only \(C\).

Only \(A\):
\begin{align*}
|A\text{ only}|
&=12000-7000-4500+500\\
&=1000.
\end{align*}

Only \(B\):
\begin{align*}
|B\text{ only}|
&=8000-7000-1000+500\\
&=500.
\end{align*}

Only \(C\):
\begin{align*}
|C\text{ only}|
&=6000-4500-1000+500\\
&=1000.
\end{align*}

Total exactly one:
\begin{align*}
1000+500+1000&=2500.
\end{align*}
So:
\begin{align*}
P(\text{exactly one})
&=\frac{2500}{30000}\\
&=\frac{1}{12}\\
&\approx 0.0833.
\end{align*}

\subsection*{Answer}
\[
\boxed{P(\text{at least one})=\frac{7}{15}\approx 0.467}
\]
\[
\boxed{P(\text{no area})=\frac{8}{15}\approx 0.533}
\]
\[
\boxed{P(\text{exactly one})=\frac{1}{12}\approx 0.0833}
\]

\subsection*{What to memorize for the exam}
For three overlapping groups, use:
\[
\text{singles} - \text{pair overlaps} + \text{triple overlap}.
\]

\subsection*{Common trap}
Do not subtract the triple overlap permanently. People in all three are included in every pair-overlap, so they need to be added back once.

\subsection*{Repeated exam pattern}
This is a new pattern so far: three-set Venn/inclusion-exclusion counting. It is related to the event-probability logic from 2026-01 Uppgift 1 and 2025-08 Uppgift 1.

\section*{Uppgift 2}

\subsection*{What this question is really asking}
Screw lengths are normally distributed:
\[
X\sim N(1.8,0.4).
\]
A screw is usable if its length is less than \(2\) cm.

We must find:
\begin{itemize}[leftmargin=*]
    \item probability one screw is usable,
    \item probability all 10 screws in a box are usable,
    \item how many screws are needed so the chance of at least one usable screw is above \(0.99\),
    \item a one-sided hypothesis test for whether the mean is now larger than \(1.8\).
\end{itemize}

\subsection*{How to recognize this type of question}
The clues are:
\begin{itemize}[leftmargin=*]
    \item It says normally distributed, so use normal standardization.
    \item Several screws in a box means multiply independent probabilities.
    \item ``At least one'' usually means use the complement ``none''.
    \item ``Greater than 1.8'' means a right-tailed hypothesis test.
\end{itemize}

\subsection*{Tiny theory needed}
To use the standard normal table, convert:
\[
z=\frac{x-\mu}{\sigma}.
\]
Here \(z\) says how many standard deviations \(x\) is from the mean.

For at least one:
\[
P(\text{at least one})=1-P(\text{none}).
\]

For a known-\(\sigma\) test:
\[
z=\frac{\bar{x}-\mu_0}{\sigma/\sqrt{n}}.
\]

\subsection*{Formula and what everything means}
\[
z=\frac{x-\mu}{\sigma}
\]
where:
\begin{itemize}[leftmargin=*]
    \item \(x\) is the boundary, here \(2\),
    \item \(\mu=1.8\) is the mean,
    \item \(\sigma=0.4\) is the standard deviation.
\end{itemize}

For all 10 usable:
\[
P(\text{all 10 usable})=p^{10},
\]
where \(p\) is the probability one screw is usable.

\subsection*{Step-by-step solution}
\textbf{Part a: One screw usable.}

Usable means:
\[
X<2.
\]
Standardize:
\begin{align*}
z
&=\frac{2-1.8}{0.4}\\
&=\frac{0.2}{0.4}\\
&=0.5.
\end{align*}
So:
\[
P(X<2)=P(Z<0.5)=\Phi(0.5)\approx 0.6915.
\]

\textbf{Part b: All 10 usable.}

Let:
\[
p=0.6915.
\]
Assuming screws are independent:
\begin{align*}
P(\text{all 10 usable})
&=p^{10}\\
&=0.6915^{10}\\
&\approx 0.0250.
\end{align*}

\textbf{Part c: At least one usable with probability above 0.99.}

At least one usable:
\[
1-P(\text{none usable}).
\]
The probability one screw is not usable is:
\[
1-p=1-0.6915=0.3085.
\]
For \(n\) screws:
\[
P(\text{none usable})=(0.3085)^n.
\]
We want:
\begin{align*}
1-(0.3085)^n&>0.99.
\end{align*}
Move terms:
\begin{align*}
(0.3085)^n&<0.01.
\end{align*}
Take logarithms:
\begin{align*}
n\ln(0.3085)&<\ln(0.01).
\end{align*}
Since \(\ln(0.3085)<0\), division flips the inequality:
\begin{align*}
n&>\frac{\ln(0.01)}{\ln(0.3085)}\\
&\approx 3.92.
\end{align*}
So the smallest whole number is:
\[
n=4.
\]

\textbf{Part d: Hypothesis test.}

We test:
\[
H_0:\mu=1.8,\qquad H_1:\mu>1.8.
\]
Given:
\[
n=10,\qquad \bar{x}=2.0,\qquad \sigma=0.4.
\]
Test statistic:
\begin{align*}
z
&=\frac{2.0-1.8}{0.4/\sqrt{10}}\\
&=\frac{0.2}{0.1265}\\
&\approx 1.58.
\end{align*}
For a right-tailed 5\% test:
\[
z_{0.95}\approx 1.645.
\]
Since:
\[
1.58<1.645,
\]
we do not reject \(H_0\).

\subsection*{Answer}
\[
\boxed{P(\text{one screw usable})\approx 0.6915}
\]
\[
\boxed{P(\text{all 10 usable})\approx 0.0250}
\]
\[
\boxed{n=4\text{ screws are needed for probability above }0.99}
\]
\[
\boxed{\text{Do not reject }H_0\text{ at 5\%.}}
\]

\subsection*{What to memorize for the exam}
For normal probabilities, standardize:
\[
z=\frac{x-\mu}{\sigma}.
\]
For at least one, use:
\[
1-P(\text{none}).
\]

\subsection*{Common trap}
In the test, do not reject just because \(\bar{x}=2.0\) is bigger than \(1.8\). It must be big enough compared with the standard error.

\subsection*{Repeated exam pattern}
The hypothesis test is the same known-\(\sigma\), one-sided mean-test pattern as 2025-08 Uppgift 6 and 2025-06 Uppgift 6.

\section*{Uppgift 3}

\subsection*{What this question is really asking}
We are given:
\[
f(x)=
\begin{cases}
ax^3, & 0<x<1,\\
0, & \text{otherwise}.
\end{cases}
\]
We need to normalize it, find the CDF, compute a probability, and find the density of:
\[
Y=\ln X.
\]

\subsection*{How to recognize this type of question}
The clues are:
\begin{itemize}[leftmargin=*]
    \item It gives a density with an unknown constant \(a\).
    \item It asks for \(F_X\), so we integrate to get the CDF.
    \item It defines a new random variable \(Y=\ln X\), so this is a transformation question.
\end{itemize}

\subsection*{Tiny theory needed}
A density must integrate to 1.

For a transformation \(Y=\ln X\), first find the CDF:
\[
F_Y(y)=P(Y\le y).
\]
Then differentiate \(F_Y(y)\) to get the density \(f_Y(y)\).

\subsection*{Formula and what everything means}
\[
F_Y(y)=P(Y\le y)=P(\ln X\le y).
\]
Since the exponential function reverses the log:
\[
\ln X\le y \quad \Longleftrightarrow \quad X\le e^y.
\]

\subsection*{Step-by-step solution}
\textbf{Part a: Find \(a\).}
\begin{align*}
\int_0^1 ax^3\,dx&=1\\
a\left[\frac{x^4}{4}\right]_0^1&=1\\
\frac{a}{4}&=1\\
a&=4.
\end{align*}

\textbf{Part b: Find \(F_X(x)\).}

For \(0<x<1\):
\begin{align*}
F_X(x)
&=\int_0^x 4t^3\,dt\\
&=[t^4]_0^x\\
&=x^4.
\end{align*}
So:
\[
F_X(x)=
\begin{cases}
0, & x\le 0,\\
x^4, & 0<x<1,\\
1, & x\ge 1.
\end{cases}
\]

\textbf{Part c: Calculate \(P(|X|<1/2)\).}

Since \(X\) only takes values between 0 and 1, \(|X|=X\). So:
\begin{align*}
P(|X|<1/2)
&=P(X<1/2)\\
&=F_X(1/2)\\
&=\left(\frac{1}{2}\right)^4\\
&=\frac{1}{16}.
\end{align*}

\textbf{Part d: Density of \(Y=\ln X\).}

Since \(0<X<1\), we know:
\[
Y=\ln X<0.
\]
So \(Y\) has support \(y<0\).

For \(y<0\):
\begin{align*}
F_Y(y)
&=P(Y\le y)\\
&=P(\ln X\le y)\\
&=P(X\le e^y)\\
&=F_X(e^y)\\
&=(e^y)^4\\
&=e^{4y}.
\end{align*}
Now differentiate:
\[
f_Y(y)=\frac{d}{dy}e^{4y}=4e^{4y}.
\]
Therefore:
\[
f_Y(y)=
\begin{cases}
4e^{4y}, & y<0,\\
0, & \text{otherwise}.
\end{cases}
\]

\subsection*{Answer}
\[
\boxed{a=4}
\]
\[
\boxed{
F_X(x)=
\begin{cases}
0, & x\le 0,\\
x^4, & 0<x<1,\\
1, & x\ge 1
\end{cases}}
\]
\[
\boxed{P(|X|<1/2)=\frac{1}{16}}
\]
\[
\boxed{
f_Y(y)=
\begin{cases}
4e^{4y}, & y<0,\\
0, & \text{otherwise}
\end{cases}}
\]

\subsection*{What to memorize for the exam}
For transformations, a beginner-safe method is:
\[
F_Y(y)=P(Y\le y),
\]
then solve the inequality back to \(X\).

\subsection*{Common trap}
Do not forget the support. Since \(0<X<1\), \(\ln X\) is always negative.

\subsection*{Repeated exam pattern}
Parts a--c repeat the continuous density pattern from 2026-01 Uppgift 2, 2025-08 Uppgift 2, and 2025-06 Uppgift 2. Part d connects to 2025-06 Uppgift 3c, where a transformed variable's distribution was found.

\section*{Uppgift 4}

\subsection*{What this question is really asking}
\(X\) is uniformly distributed on \((-1,1)\). That means all values between \(-1\) and \(1\) are equally likely.

\(Y\)'s distribution is symmetric around 0. That means:
\[
P(Y\le 0)=\frac{1}{2}.
\]
Also \(X\) and \(Y\) are independent.

\subsection*{How to recognize this type of question}
The clues are:
\begin{itemize}[leftmargin=*]
    \item ``likformigt fördelad'' means uniform distribution.
    \item ``symmetrisk kring origo'' means half the probability is below 0 and half above 0.
    \item \(\max(X,Y)\le 0\) means both \(X\le 0\) and \(Y\le 0\).
\end{itemize}

\subsection*{Tiny theory needed}
For a uniform distribution on \((a,b)\), probability is interval length divided by total length:
\[
P(c<X<d)=\frac{d-c}{b-a}.
\]
For \(X\sim U(-1,1)\):
\[
E(X)=0,\qquad V(X)=\frac{(1-(-1))^2}{12}=\frac{1}{3}.
\]

\subsection*{Formula and what everything means}
For \(W=2X-1\):
\[
F_W(w)=P(W\le w)=P(2X-1\le w).
\]
Then solve:
\[
X\le \frac{w+1}{2}.
\]

\subsection*{Step-by-step solution}
\textbf{Part a.}
\begin{align*}
P(-0.1<X<0.2)
&=\frac{0.2-(-0.1)}{1-(-1)}\\
&=\frac{0.3}{2}\\
&=0.15.
\end{align*}

\textbf{Part b.}
\begin{align*}
E(2X-1)
&=2E(X)-1\\
&=2\cdot 0-1\\
&=-1.
\end{align*}
Variance:
\begin{align*}
V(2X-1)
&=2^2V(X)\\
&=4\cdot \frac{1}{3}\\
&=\frac{4}{3}.
\end{align*}

\textbf{Part c.}

Let:
\[
W=2X-1.
\]
Since \(-1<X<1\), the new variable satisfies:
\[
-3<W<1.
\]
For \(-3<w<1\):
\begin{align*}
F_W(w)
&=P(W\le w)\\
&=P(2X-1\le w)\\
&=P\left(X\le \frac{w+1}{2}\right).
\end{align*}
The CDF of \(X\sim U(-1,1)\) is:
\[
F_X(x)=\frac{x+1}{2},\qquad -1<x<1.
\]
Insert \(x=(w+1)/2\):
\begin{align*}
F_W(w)
&=\frac{\frac{w+1}{2}+1}{2}\\
&=\frac{w+3}{4}.
\end{align*}
So:
\[
F_W(w)=
\begin{cases}
0, & w\le -3,\\
\dfrac{w+3}{4}, & -3<w<1,\\
1, & w\ge 1.
\end{cases}
\]

\textbf{Part d.}

\[
\max(X,Y)\le 0
\]
means:
\[
X\le 0 \text{ and } Y\le 0.
\]
By independence:
\begin{align*}
P(\max(X,Y)\le 0)
&=P(X\le 0)P(Y\le 0)\\
&=\frac{1}{2}\cdot \frac{1}{2}\\
&=\frac{1}{4}.
\end{align*}

\subsection*{Answer}
\[
\boxed{P(-0.1<X<0.2)=0.15}
\]
\[
\boxed{E(2X-1)=-1,\qquad V(2X-1)=\frac{4}{3}}
\]
\[
\boxed{
F_{2X-1}(w)=
\begin{cases}
0, & w\le -3,\\
\frac{w+3}{4}, & -3<w<1,\\
1, & w\ge 1
\end{cases}}
\]
\[
\boxed{P(\max(X,Y)\le 0)=\frac{1}{4}}
\]

\subsection*{What to memorize for the exam}
Uniform probability is length over total length. For max questions:
\[
\max(X,Y)\le c \iff X\le c \text{ and } Y\le c.
\]

\subsection*{Common trap}
For variance, the constant \(-1\) does not matter, but the coefficient \(2\) gets squared:
\[
V(2X-1)=4V(X).
\]

\subsection*{Repeated exam pattern}
Part c repeats the transformed-variable CDF pattern from 2025-06 Uppgift 3c. Part d repeats the maximum pattern from 2025-06 Uppgift 3d and 2025-08 Uppgift 3b.

\section*{Uppgift 5}

\subsection*{What this question is really asking}
We have a density depending on \(\theta\):
\[
f_X(x)=\theta(1+x)^{-1-\theta},\qquad x\ge 0.
\]
We know \(\theta\) is only allowed to be:
\[
2,\quad 3,\quad 4.
\]
The sample values are:
\[
0.2,\quad 0.8.
\]
We must choose which \(\theta\) makes the sample most likely.

\subsection*{How to recognize this type of question}
The clue is ``ML-skattningen''. Usually we derive by differentiating, but here \(\theta\) can only be 2, 3, or 4. So the easiest method is to calculate the likelihood for each and choose the largest.

\subsection*{Tiny theory needed}
Likelihood means:
\[
L(\theta)=f(0.2)f(0.8).
\]
The ML estimate is the \(\theta\) that gives the largest \(L(\theta)\).

\subsection*{Formula and what everything means}
\[
L(\theta)
=\theta(1+0.2)^{-1-\theta}\cdot \theta(1+0.8)^{-1-\theta}.
\]
This becomes:
\[
L(\theta)=\theta^2(1.2)^{-1-\theta}(1.8)^{-1-\theta}.
\]

\subsection*{Step-by-step solution}
Since:
\[
1.2\cdot 1.8=2.16,
\]
we can write:
\[
L(\theta)=\theta^2(2.16)^{-1-\theta}.
\]

Now test the allowed values.

For \(\theta=2\):
\[
L(2)=2^2(2.16)^{-3}\approx 0.3969.
\]

For \(\theta=3\):
\[
L(3)=3^2(2.16)^{-4}\approx 0.4135.
\]

For \(\theta=4\):
\[
L(4)=4^2(2.16)^{-5}\approx 0.3403.
\]

The largest likelihood is:
\[
L(3)\approx 0.4135.
\]
Therefore:
\[
\hat{\theta}_{ML}=3.
\]

\subsection*{Answer}
\[
\boxed{\hat{\theta}_{ML}=3}
\]

\subsection*{What to memorize for the exam}
If the parameter is only allowed to take a few values, you can simply calculate the likelihood for each value and choose the biggest.

\subsection*{Common trap}
Do not give the continuous maximizer if the question says \(\theta\) is either 2, 3, or 4. The answer must be one of those three.

\subsection*{Repeated exam pattern}
This is the same ML pattern as 2026-01 Uppgift 5, 2025-08 Uppgift 4, and 2025-06 Uppgift 4, but with a discrete set of possible parameter values.

\section*{Uppgift 6}

\subsection*{What this question is really asking}
We have two groups of students and three possible grades:
\[
F,\quad G,\quad VG.
\]
We test whether the grade distribution is the same in both groups.

This is a chi-square test for homogeneity.

\subsection*{How to recognize this type of question}
The clues are:
\begin{itemize}[leftmargin=*]
    \item Data are in a table of counts.
    \item There are two groups and several categories.
    \item The question asks whether the distributions are the same.
\end{itemize}

\subsection*{Tiny theory needed}
The chi-square statistic is:
\[
\chi^2=\sum \frac{(O-E)^2}{E}.
\]
Here:
\begin{itemize}[leftmargin=*]
    \item \(O\) is observed count,
    \item \(E\) is expected count if the distributions were the same.
\end{itemize}

Expected count:
\[
E=\frac{\text{row total}\cdot \text{column total}}{\text{grand total}}.
\]

\subsection*{Formula and what everything means}
Degrees of freedom:
\[
(r-1)(c-1).
\]
Here:
\begin{itemize}[leftmargin=*]
    \item \(r=2\) rows,
    \item \(c=3\) grade categories.
\end{itemize}
So:
\[
(2-1)(3-1)=2.
\]

\subsection*{Step-by-step solution}
Observed table:
\[
\begin{array}{c|ccc|c}
&F&G&VG&\text{row total}\\ \hline
\text{group 1}&15&10&5&30\\
\text{group 2}&23&27&10&60\\ \hline
\text{column total}&38&37&15&90
\end{array}
\]

Hypotheses:
\[
H_0:\text{grade distribution is the same in both groups}
\]
\[
H_1:\text{grade distribution is not the same in both groups}.
\]

Expected counts for group 1:
\[
E_{1,F}=\frac{30\cdot 38}{90}=12.667,
\]
\[
E_{1,G}=\frac{30\cdot 37}{90}=12.333,
\]
\[
E_{1,VG}=\frac{30\cdot 15}{90}=5.
\]

Expected counts for group 2:
\[
E_{2,F}=\frac{60\cdot 38}{90}=25.333,
\]
\[
E_{2,G}=\frac{60\cdot 37}{90}=24.667,
\]
\[
E_{2,VG}=\frac{60\cdot 15}{90}=10.
\]

Now compute the statistic:
\begin{align*}
\chi^2
&=\frac{(15-12.667)^2}{12.667}
+\frac{(10-12.333)^2}{12.333}
+\frac{(5-5)^2}{5}\\
&\quad+\frac{(23-25.333)^2}{25.333}
+\frac{(27-24.667)^2}{24.667}
+\frac{(10-10)^2}{10}\\
&\approx 1.307.
\end{align*}

At 5\% significance with 2 degrees of freedom:
\[
\chi^2_{0.95,2}\approx 5.991.
\]
Since:
\[
1.307<5.991,
\]
we do not reject \(H_0\).

\subsection*{Answer}
\[
\boxed{\chi^2\approx 1.307}
\]
\[
\boxed{\text{Do not reject }H_0.}
\]
There is not enough evidence at the 5\% level to say that the grade distribution differs between the two groups.

\subsection*{What to memorize for the exam}
For a table of counts:
\[
E=\frac{\text{row total}\cdot \text{column total}}{\text{grand total}},
\qquad
\chi^2=\sum\frac{(O-E)^2}{E}.
\]

\subsection*{Common trap}
Do not compare percentages by eye only. The chi-square test measures whether the differences are large relative to the expected counts.

\subsection*{Repeated exam pattern}
This is a new pattern in the processed folders: chi-square test for homogeneity in a contingency table.

\section*{Mini-cheatsheet from this exam}

\subsection*{Venn diagrams}
\[
|A\cup B\cup C|=|A|+|B|+|C|-|AB|-|AC|-|BC|+|ABC|.
\]

\subsection*{Normal distribution}
\[
z=\frac{x-\mu}{\sigma}.
\]
Known-\(\sigma\) test:
\[
z=\frac{\bar{x}-\mu_0}{\sigma/\sqrt{n}}.
\]

\subsection*{At least one}
\[
P(\text{at least one})=1-P(\text{none}).
\]

\subsection*{Density and transformations}
\[
\int f(x)\,dx=1,\qquad F_X(x)=\int_{-\infty}^{x}f(t)\,dt.
\]
For transformations:
\[
F_Y(y)=P(Y\le y).
\]

\subsection*{Uniform distribution}
For \(X\sim U(a,b)\):
\[
P(c<X<d)=\frac{d-c}{b-a},\qquad V(X)=\frac{(b-a)^2}{12}.
\]

\subsection*{Maximum}
\[
\max(X,Y)\le c \iff X\le c \text{ and } Y\le c.
\]

\subsection*{Maximum likelihood}
If only a few parameter values are allowed, calculate \(L\) for each and choose the largest.

\subsection*{Chi-square table test}
\[
E=\frac{\text{row total}\cdot \text{column total}}{\text{grand total}},
\qquad
\chi^2=\sum\frac{(O-E)^2}{E}.
\]

\end{document}
